\section{A regularized one-sided Euler commutator inside the critical strip}
\label{sec:regularized-strip}

This section constructs actual closed multiplier operators and an
actual trace-class commutator on almost every one-sided weight inside
the critical strip.  The test and the derivative of the cutoff have
compact support and finite Gevrey order.  The trace is obtained from
the operator kernel before a contour shift is made.  The shift then
exhibits the pole and the zeros crossed by this regularization.
No bounded M\"obius inverse on a matching narrow Hilbert space is
asserted, and the result is not stated for every width or for every
compact smooth test.

For a finite set $S$ of rational primes write
\[
 E_S(s)=\prod_{p\in S}(1-p^{-s}),\qquad
 \zeta_S(s)=E_S(s)\zeta(s).
 \tag{MA-RS.1}
\]
The subscript removes the primes in $S$ from the Euler product.
Let $\rho=\beta+i\gamma$ run through the distinct nontrivial
zeros of $\zeta$, with multiplicities $m_\rho$.

\paragraph{MA-RS-001 (generic vertical lines and the inverse multiplier).}
There is a full-measure subset
$\mathscr G\subset(1/2,1)$ such that for every
$\sigma\in\mathscr G$ there is $c_\sigma>0$ with
\[
       |\sigma-\beta|
          \ge c_\sigma(2+|\gamma|)^{-2}
                    \quad\hbox{for every nontrivial zero }\rho.
 \tag{MA-RS.2}
\]
For this $\sigma$ and every integer $r\ge0$ there is
$C_{\sigma,S,r}>0$ such that
\[
 \left|\frac{d^r}{dt^r}\frac1{\zeta_S(\sigma+it)}\right|
       \le \exp\!\left(C_{\sigma,S,r}
                            \log^2(2+|t|)\right)
                    \qquad(t\in\mathbb R).
 \tag{MA-RS.3}
\]
The functions $\zeta_S(\sigma+it)$ and all of their
$t$ derivatives have polynomial growth.  These assertions
are unconditional.

\begin{proof}
The unconditional local zero count used in NT-Z10 and
MT-M18 is $O(\log(2+|T|))$ in each fixed-height window,
with multiplicities.  Consequently
\[
                  \sum_\rho (2+|\gamma|)^{-2}<\infty.
 \tag{MA-RS.4}
\]
The intervals centered at $\beta$ of radius
$(2+|\gamma|)^{-2}$ have summable lengths.
The first Borel--Cantelli lemma says that almost every
$\sigma$ belongs to only finitely many of them.
Remove also the countable set of exact real parts $\beta$.
For each remaining $\sigma$, the finitely many exceptional
distances have a positive minimum after division by their
respective radii.  Decreasing that minimum to at most one
proves (MA-RS.2).

We give the needed minimum-modulus argument, because a
zero-free vertical line alone would not imply (MA-RS.3).
At large $|t|$, center a disk at $2+it$ and choose a
zero-free boundary radius $R$ in a fixed sufficiently large
interval, for example $12<R<13$.
The disk contains $M=O(\log(2+|t|))$ zeros and no pole or
trivial zero of $\zeta$.  Divide them out with the finite
disk Blaschke product $B$, with multiplicities, and write
$\zeta=G B$.  On the boundary $|B|=1$.
The fixed-strip polynomial upper bound gives
$|G|\le C(2+|t|)^C$ inside the disk.  At its center,
$|G(2+it)|\ge|\zeta(2+it)|\ge\zeta(2)^{-1}$.
Harnack's inequality, applied to
$\log(C(2+|t|)^C)-\log|G|$, therefore gives
\[
          |G(w)|\ge c(2+|t|)^{-C}
 \tag{MA-RS.5}
\]
on the inner half disk.  This is the same finite-disk
argument used in the proof of MT-M18 (M18.8).

Take $w=\sigma+it$.  In the normalized disk coordinate
it lies in the inner quarter disk, so each Blaschke
factor corresponding to $\rho$ is at least
$|w-\rho|/(2R)$.  For every zero in the disk,
(MA-RS.2) and $|\gamma-t|<13$ give
\[
       |w-\rho|\ge c'_\sigma(2+|t|)^{-2}.
 \tag{MA-RS.6}
\]
Multiplying at most $C\log(2+|t|)$ such lower bounds and
then using (MA-RS.5) yields
\[
        |\zeta(\sigma+it)|^{-1}
              \le \exp(C_\sigma\log^2(2+|t|)).
 \tag{MA-RS.7}
\]
The same argument works uniformly when $w$ ranges over
a disk of radius $c''_\sigma(2+|t|)^{-2}$ about
$\sigma+it$, after halving the constant in (MA-RS.6).
That disk is zero free.  Cauchy's estimate for $1/\zeta$
on it proves (MA-RS.3) for $S=\varnothing$: the extra
factor $r!c_\sigma^{-r}(2+|t|)^{2r}$ is absorbed into
the right side.  A compact range of $t$ is covered by
enlarging the constants.

On a small closed strip about $\sigma>1/2$, all factors
$1-p^{-s}$, $p\in S$, are bounded away from zero, and
their inverses and derivatives are bounded.
This proves (MA-RS.3) for finite $S$.
The fixed-strip upper bound for $\zeta$, with Cauchy
estimates for its derivatives, proves the remaining
polynomial-growth assertion.  The pole at $1$ is
avoided by the fixed line and causes no issue on the
compact low-height portion.
\end{proof}

\subsection{The closed operators and their dense common core}

Fix $\sigma\in\mathscr G$ and put $a=\sigma-1/2$.
Use
\[
 \begin{aligned}
 \mathcal H_a&=L^2(\mathbb R,e^{2ax}\,dx),&
 W_au(x)&=e^{ax}u(x),\\
 \mathcal F_+v(t)&=(2\pi)^{-1/2}
                       \int_{\mathbb R}e^{itx}v(x)\,dx,&
 V_a&=\mathcal F_+W_a,\\
 \mathscr D_a&=V_a^{-1}C_c^\infty(\mathbb R),&
 A(t)&=\zeta_S(\sigma+it),\qquad j(t)=A(t)^{-1}.
 \end{aligned}
 \tag{MA-RS.8}
\]
Thus $V_a$ is unitary and $\mathscr D_a$ is the dense
core with compact support in the weighted Fourier variable.
Define
\[
 \begin{aligned}
 Z_{a,S}&=V_a^{-1}M_A V_a,&
 \mathcal D(M_A)&=\{v\in L^2:Av\in L^2\},\\
 J_{a,S}&=V_a^{-1}M_j V_a,&
 \mathcal D(M_j)&=\{v\in L^2:jv\in L^2\}.
 \end{aligned}
 \tag{MA-RS.9}
\]
These are closed, densely defined inverse operators, with
their maximal multiplier domains.  The equality
$M_A M_j=1$ holds on $\mathcal D(M_j)$, and
$M_j M_A=1$ on $\mathcal D(M_A)$.
Both have $C_c^\infty$ as a graph core in the Fourier
picture, by cutoff and smooth approximation in their
weighted graph norms.  Neither inverse identity says
that both operators are bounded on all of $\mathcal H_a$.

This continuation of $Z$ is also an explicit regularization
of its Euler sums.  For every positive integer $N$ and
$\operatorname{Re}s>0$, $s\ne1$, Abel summation gives
\[
 \zeta(s)=\sum_{n\le N}n^{-s}+\frac{N^{1-s}}{s-1}
           -s\int_N^\infty\{x\}x^{-s-1}\,dx .
 \tag{MA-RS.10}
\]
The remainder has modulus at most
$|s|N^{-\operatorname{Re}s}/\operatorname{Re}s$.
In particular the bounded multipliers
\[
        E_S(s)\left(\sum_{n\le N}n^{-s}
                              +\frac{N^{1-s}}{s-1}\right),
                  \qquad s=\sigma+it,
 \tag{MA-RS.11}
\]
converge to $A(t)$ on each compact frequency interval,
and their operators converge to $Z_{a,S}$ on
$\mathscr D_a$.  The closed graph closure gives (MA-RS.9).
Formula (MA-RS.10) follows first in $\operatorname{Re}s>1$
by writing $\lfloor x\rfloor=x-\{x\}$ in the tail
Stieltjes integral, and then continues to
$\operatorname{Re}s>0$ because the displayed integral
converges there.  Its exact published reference is
DLMF~25.2.8, equivalently Apostol (1976),
Theorem~12.21 and equation~(25), p.~269.
For $\operatorname{Re}s>1$ the counterterm tends to zero,
and this agrees with the operator-norm Euler construction
of MA-MEY-001.  Inside the critical strip the counterterm
is part of the definition, not an omitted convergence claim.

For completeness the inverse can be approximated on its
actual domain by the bounded multipliers
\[
                j_\varepsilon(t)
                  =\frac{\overline{A(t)}}{|A(t)|^2+\varepsilon}.
 \tag{MA-RS.12}
\]
They converge to $j$ in $L^2$ after multiplication by any
vector in $\mathcal D(M_j)$, since $|j_\varepsilon|\le|j|$.
This is a domain approximation; no uniform bound for their
operator norms as $\varepsilon\downarrow0$ is asserted.

\subsection{Gevrey decay and an actual trace-class kernel}

A compactly supported function has Gevrey order $q>1$
if its derivatives satisfy
$\|f^{(N)}\|_\infty\le C R^N(N!)^q$ for every $N\ge0$.
Choose $g$ of compact support and finite Gevrey order,
and a bounded smooth cutoff $\phi$ such that
\[
 \phi(-\infty)=1,\qquad \phi(+\infty)=0,\qquad
                 \phi'\text{ has compact support and finite Gevrey order}.
 \tag{MA-RS.13}
\]
Increase $q>1$ so that one common order works for $g$ and
$\phi'$.  Compact Gevrey bumps yield such cutoffs by
integration; a symmetric bump gives
$\phi(x)+\phi(-x)=1$ if that additional convention is desired.
Put
\[
 G(z)=\mathcal Mg(z)=\int e^{zx}g(x)\,dx,\qquad
 m(t)=G(a+it),\qquad
 p(\xi)=\int e^{i\xi x}\phi'(x)\,dx .
 \tag{MA-RS.14}
\]
The transforms here, unlike $\mathcal F_+$, are not
normalized by $(2\pi)^{-1/2}$.

\paragraph{MA-RS-002 (the regularized commutator is defined on its core).}
On $\mathscr D_a$ the composition
\[
           K_{a,S,\phi}(g)
              =C_gZ_{a,S}[J_{a,S},M_\phi]
 \tag{MA-RS.15}
\]
is well defined.  It extends uniquely to a trace-class
operator on $\mathcal H_a$.  In the unitary $V_a$ picture
its kernel with respect to $du$ is the Schwartz function
\[
 \mathcal K(t,u)=\frac{i}{2\pi}\,
          m(t)A(t)\frac{j(t)-j(u)}{t-u}\,p(t-u),
 \tag{MA-RS.16}
\]
with its smooth divided-difference value at $t=u$.
Its actual operator trace is
\[
       \operatorname{Tr}K_{a,S,\phi}(g)
        =-\frac1{2\pi}\int_{\mathbb R}
               G(a+it)\frac{\zeta_S'}{\zeta_S}(\sigma+it)\,dt .
 \tag{MA-RS.17}
\]
The integral is absolutely convergent.  The trace-norm
estimates are uniform when test supports and Gevrey
derivative constants are fixed.

\begin{proof}
We first give the decay estimate used both for domains
and for the kernel.  Repeated integration by parts gives
for a compact Gevrey function $f$
\[
       |\widehat f_+(t)|
          \le C R^N(N!)^q|t|^{-N}\qquad(N\ge0).
 \tag{MA-RS.18}
\]
Choose $N$ to be a fixed positive multiple of
$|t|^{1/q}$, rounded down, so that
$R N^q/|t|\le1/2$, and use $N!\le N^N$.
This proves $|\widehat f_+(t)|\le C'e^{-c|t|^{1/q}}$.
Multiplying $f$ by any fixed power of $x$ or by $e^{ax}$
preserves its Gevrey order on the fixed compact support,
by Leibniz's formula.  Hence for every fixed $r\ge0$,
\[
          |m^{(r)}(t)|+|p^{(r)}(t)|
                 \le C_r e^{-c_r|t|^{1/q}}.
 \tag{MA-RS.19}
\]
These estimates are uniform for the bounded Gevrey
families specified in the statement.

Convolution satisfies $V_a C_gV_a^{-1}=M_m$.
Let $v\in C_c^\infty$ be a weighted Fourier core vector,
and put $f=V_a^{-1}v\in\mathscr D_a$.
Then $jv\in C_c^\infty$.
For a bounded smooth $\phi$ with compactly supported
derivative, $\phi\,\mathcal F_+^{-1}v$ is Schwartz,
because $\mathcal F_+^{-1}v$ is Schwartz and all
derivatives of $\phi$ are bounded.
Outside a sufficiently large interval containing
$\operatorname{supp}v$, its Fourier transform is the
ordinary integral
\[
    \frac{i}{2\pi}\int
                 \frac{p(t-u)}{t-u}\,v(u)\,du .
 \tag{MA-RS.20}
\]
Indeed the distribution identity
$\widehat{\phi'}_+(\xi)=-i\xi\widehat\phi_+(\xi)$
gives $\widehat\phi_+(\xi)=i p(\xi)/\xi$
away from zero.  The same formula and its derivatives
show Gevrey exponential decay in that exterior region,
by (MA-RS.19).  The assertion applies equally with
$v$ replaced by $jv$.

It follows from (MA-RS.3) that $M_\phi$ applied to a
core vector lies in $\mathcal D(J_{a,S})$:
the tails from (MA-RS.20) dominate
$\exp(C\log^2(2+|t|))$.
Moreover $J_{a,S}M_\phi f$ lies in
$\mathcal D(Z_{a,S})$, and $Z_{a,S}J_{a,S}M_\phi f=M_\phi f$.
The second composition $Z_{a,S}M_\phi J_{a,S}f$
is also defined, because the Fourier transform after
$M_\phi J_{a,S}$ has those same tails and $A$
has polynomial growth.  This verifies all domains in
(MA-RS.15), including the application of $Z$ to the
commutator.  On this core it also proves
\[
       K_{a,S,\phi}(g)
          =C_g(M_\phi-Z_{a,S}M_\phi J_{a,S}).
 \tag{MA-RS.21}
\]

In Fourier variables multiplication by $\phi$ has
distribution kernel $(2\pi)^{-1}\widehat\phi_+(t-u)$.
The commutator multiplies this by $j(t)-j(u)$.
Since this difference vanishes at $t=u$, the possible
delta term in $\widehat\phi_+$ disappears, and the
identity $\xi\widehat\phi_+=i p$ gives precisely
(MA-RS.16).  This use of a distribution kernel is
initially on compact-frequency core vectors, where
all its products are defined; the resulting kernel
will now be seen to extend to all of $L^2$.

Write
\[
       \frac{j(t)-j(u)}{t-u}
           =\int_0^1 j'(u+\theta(t-u))\,d\theta .
 \tag{MA-RS.22}
\]
Every derivative of this function is bounded by
$\exp(C\log^2(2+|t|+|u|))$, with a derivative-dependent
constant, by (MA-RS.3).
All derivatives of $A$ have polynomial growth.
Combining (MA-RS.19) with the product rule gives,
for each fixed pair of derivative orders, an upper
bound of the form
\[
 C\exp\!\left(
     -c|t|^{1/q}-c|t-u|^{1/q}
             +C'\log^2(2+|t|+|u|)\right).
 \tag{MA-RS.23}
\]
The logarithmic term is absorbed by half the two
negative Gevrey powers; use
$|u|\le|t|+|t-u|$.
The same bound therefore dominates every polynomial
weight in $(t,u)$.  This proves
$\mathcal K\in\mathcal S(\mathbb R^2)$.

A Schwartz kernel defines a trace-class operator,
and its trace is its diagonal integral.  One explicit
justification is to expand in normalized Hermite
functions in the two variables.  Applying an arbitrary
power of the harmonic oscillator in each variable
bounds the coefficients by
$C_N(1+r)^{-N}(1+s)^{-N}$.
For $N>1$ their absolute sum is finite, so the
rank-one series converges in trace norm.
Higher powers give convergence of the kernel series
in Schwartz seminorms; integration of the diagonal
then gives the same trace.  This also proves the
claimed uniform trace-norm bounds from finitely many
Schwartz seminorms of (MA-RS.16).
Agreement on the dense core proves uniqueness of
the bounded extension.

Finally $p(0)=\int\phi'=-1$, and
\[
             j'(t)=-i\frac{\zeta_S'(\sigma+it)}
                                  {\zeta_S(\sigma+it)^2}.
 \tag{MA-RS.24}
\]
Thus the diagonal of (MA-RS.16) is
$-(2\pi)^{-1}G(a+it)(\zeta_S'/\zeta_S)(\sigma+it)$.
It is integrable by (MA-RS.3) and (MA-RS.19).
This proves (MA-RS.17), including its sign.
\end{proof}

\subsection{The residues carried by this operator regularization}

For compact $g$ put
\[
 \begin{aligned}
 P_{\overline S}^{+}(g)&=
       \sum_{\substack{p\notin S\\k\ge1}}
          (\log p)p^{-k/2}g(k\log p),\\
 B_{\overline S}(g)&=
          P_{\overline S}^{+}(g)+P_{\overline S}^{+}(g^\vee),
       &g^\vee(x)&=g(-x),\\
 P(g)&=G(1/2)+G(-1/2).
 \end{aligned}
 \tag{MA-RS.25}
\]
Both prime sums are finite.

\paragraph{MA-RS-003 (crossed zeros are an actual commutator correction).}
For the operators and tests just constructed,
\[
 \operatorname{Tr}K_{a,S,\phi}(g)
       =P_{\overline S}^{+}(g)-G(1/2)
                +\sum_{\operatorname{Re}\rho>\sigma}
                             m_\rho G(\rho-1/2).
 \tag{MA-RS.26}
\]
The zero sum is absolutely convergent.
Let $R:\mathcal H_a\to\mathcal H_{-a}$ be reflection,
$Ru(x)=u(-x)$.  The paired diagonal regularization
on $\mathcal H_a\oplus\mathcal H_{-a}$ is the actual
trace-class operator
\[
        \mathbb K_{a,S,\phi}(g)
             =K_{a,S,\phi}(g)\oplus
                       R K_{a,S,\phi}(g^\vee)R .
 \tag{MA-RS.27}
\]
It satisfies
\[
 \boxed{\ 
 \operatorname{Tr}\mathbb K_{a,S,\phi}(g)
    =B_{\overline S}(g)-P(g)
       +\sum_{|\operatorname{Re}\rho-1/2|>a}
                         m_\rho G(\rho-1/2).\ }
 \tag{MA-RS.28}
\]
There are no boundary zeros at this generic width.
Thus the regularized diagonal expression retains precisely
the contributions crossed in moving the defining Mellin
lines; its trace is not the unchanged wide-strip prime sum.

\begin{proof}
Rewrite (MA-RS.17) as
\[
    -\frac1{2\pi i}\int_{\sigma-i\infty}^{\sigma+i\infty}
           G(s-1/2)\frac{\zeta_S'}{\zeta_S}(s)\,ds .
 \tag{MA-RS.29}
\]
Fix $c>1$.  On the line $\operatorname{Re}s=c$ the
absolutely convergent Euler logarithmic derivative is
\[
          -\frac{\zeta_S'}{\zeta_S}(s)
           =\sum_{\substack{p\notin S\\k\ge1}}
                              (\log p)p^{-ks}.
 \tag{MA-RS.30}
\]
Interchange this absolute series with the integrable
Mellin test.  Fourier inversion on the vertical line
gives
\[
 -\frac1{2\pi i}\int_{c-i\infty}^{c+i\infty}
        G(s-1/2)\frac{\zeta_S'}{\zeta_S}(s)\,ds
                      =P_{\overline S}^{+}(g).
 \tag{MA-RS.31}
\]
This also follows from the already established actual
wide-strip commutator trace MA-MEY-001.

Use the selected upper and lower horizontal ordinates
of MT-M18 (M18.7)--(M18.8).
Its minimum-modulus proof, together with Stirling's
upper and lower bounds for $\Xi/\zeta$, makes
$|\zeta'/\zeta|$ polynomially bounded on the horizontal
segments $\sigma\le\operatorname{Re}s\le c$ at these
heights.  The factors $E_S$ and their inverses are
uniformly bounded there since $\sigma>1/2$.
The Gevrey estimate for $G(s-1/2)$ is uniform on this
fixed real interval.  Therefore both horizontal
integrals tend to zero.  The vertical integrals
converge absolutely by the same estimates as in
MA-RS-002.

The residue of the integrand in (MA-RS.29) before
its outside minus sign is $m_\rho G(\rho-1/2)$
at a nontrivial zero and $-G(1/2)$ at the pole $s=1$.
The extra zeros of $E_S$ have real part zero and
are not crossed.  The counterclockwise rectangle
has the $c$ edge oriented upwards and the $\sigma$
edge downwards.  Thus the upward $\sigma$ integral
equals the upward $c$ integral minus $2\pi i$
times the residue sum.  Multiplying by
$-(2\pi i)^{-1}$ gives (MA-RS.26).
The local zero count and uniform vertical decay
of $G$ prove absolute convergence of the resulting
zero sum, so the chosen horizontal subsequence
does not change its value.

Reflection is unitary between the two weighted
spaces, and $RC_gR=C_{g^\vee}$.
Thus (MA-RS.27) is the smoothed pair of diagonal
expressions $D$ and $RDR$, where
$D=Z_{a,S}[J_{a,S},M_\phi]$ is initially defined
on its core.  It is not a definition by its trace.
Apply (MA-RS.26) to $g$ and $g^\vee$ and add.
Their Mellin transforms satisfy
$\mathcal M(g^\vee)(z)=G(-z)$.
The nontrivial zeta divisor is invariant, with
multiplicity, under $\rho\mapsto1-\rho$.
Consequently the second right-hand zero sum becomes
the sum over $\operatorname{Re}\rho<1-\sigma$.
These two disjoint sets are exactly
$|\operatorname{Re}\rho-1/2|>a$.
The two pole terms sum to $-P(g)$.
Equation (MA-RS.2), and its reflected version from
the same zero symmetry, exclude zeros on either
boundary.  This proves (MA-RS.28).
\end{proof}

\paragraph{Compact cutoffs and the limits of the result.}
If $\phi$ in MA-RS-002 is replaced by a compactly supported
Gevrey function, the same kernel and domain proof applies,
but $p(0)=0$, so the commutator trace is zero.
Whenever (MA-RS.26) is nonzero, no family of these compact
cutoff commutators can converge to the one-sided operator
in trace norm.  This conclusion uses only the trace
inequality and does not assert strong convergence for
the unbounded inverse compositions.

The theorem covers $\sigma\in\mathscr G$, a full-measure
set, and compact tests and cutoff derivatives of finite
Gevrey order.  A zero-free line not satisfying the
quantitative estimate has not been covered.  Arbitrary
$C_c^\infty$ Fourier decay need not dominate the inverse
growth estimate proved here, so passage to all compact
smooth tests has not been justified by this argument.
At a line containing a zero, the multiplier inverse
and boundary-domain issue must be treated separately.

There is a distinction between extending the scalar functional
and extending the operator.  The disk argument in MA-RS-001
also makes $\zeta_S'/\zeta_S$ polynomially bounded on the
chosen line.  Indeed $G'/G$ is polynomially bounded by
Cauchy's estimate and (MA-RS.5), while the logarithmic
derivative of the finite Blaschke product is bounded by
$C(2+|t|)^2\log(2+|t|)$ using (MA-RS.6).
Thus the integral in (MA-RS.17), considered only as a
scalar distribution, extends continuously to every
$C_c^\infty$ test and has the same residue identity.
This does not improve the available bounds for the
two-variable divided difference of $1/\zeta_S$, and
does not prove a trace-class operator extension for
those arbitrary smooth tests.

The paired operator in (MA-RS.27) realizes the regularized
Euler diagonal expression.  It is not a proof that all
of Meyer's common-space projection maps, with their
off-diagonal terms, extend to this critical-strip
completion.  Such an identification would require its
own domain and intertwining proof.  The residue correction
in (MA-RS.28) is already forced for this concrete
regularization; deleting it would change the trace of
the actual constructed operator.

In the requested geometric interpretation, the crossed
divisor is the spectral boundary contribution to a change
of completion.  Its occurrence is compatible with the
actual partial quotients of MA-HY and MA-WP and with
MA-HM-003's obstruction to recovering discarded blocks.
The nuclear trace viewpoint has the inherited
BR-TR-01/02 and F-SM-11/12 placement at cote~1,
scans~079--080 and 111--112.  This section does not infer
a positive Grothendieck-group realization or a
zero-location theorem from that interpretation.

\paragraph{Exact source and proof scope.}
The arithmetic inputs are the unconditional fixed-strip
upper bound and local zero count of NT-Z10, together with
the finite-disk Harnack argument and selected-height
minimum modulus in MT-M18 (M18.5)--(M18.9).
MA-RS-001 supplies the new generic-line version in full.
The explicit regularization formula is DLMF~25.2.8,
\url{https://dlmf.nist.gov/25.2.E8}, sourced there to
Apostol, T. M. (1976), \emph{Introduction to analytic
number theory}, Theorem~12.21 and equation~(25), p.~269.
The Fourier sign, core-domain argument, Gevrey estimate,
Schwartz-kernel estimate and contour correction are all
proved in this section.  MA-MEY-001--003 supplies the
already checked wide-strip operator and source
diagonal conventions; its zero-location-free source is
Meyer (2005), Lemmas~5.3--5.6 and Theorem~5.8, author
pp.~39--44.  No A, M, or new outside preprint is used.
